\scanpage{93}
% archive sheet number: 92
% block text | line 1, cancelled inset
\deleted{\unsure{Caractéristique}. Si on a}
% block diagram | line 2, cancelled inset
\diagram{\begin{tikzcd}[ampersand replacement=\&]
\deleted{0}\arrow[r]\&\deleted{M'}\arrow[r,dash]\&\deleted{M}\arrow[r,dash]\&\deleted{M''}\arrow[r]\&\deleted{0}
\end{tikzcd}}
% block text | line 3, cancelled inset
\deleted{et si \unsure{on a} \illegible{2} \unsure{propriété|injectivité} \unsure{pour}}
% block text | lines 1-8, cancelled surrounding block
\deleted{\unsure{provient} du fait que, si $\unsure{\delta|d}=[\Gamma:\Lambda]$, la multiplication par $\delta$ est une bijection dans $A$ (car \unsure{ça l'est} dans les $m^iA/m^{i+1}A$). Par suite, on \deleted{\unsure{peut trouver}} a un $v\in A$ tel que}
% block display | line 9, cancelled surrounding block
\[
\deleted{\unsure{\gamma|\sigma}v-v=ug(\gamma)\qquad\text{pour }\underline{\forall}\gamma\in\Gamma}
\]
% block text | line 10
Soit $v\in A^\Gamma$ tel que $\unsure{\varepsilon|\epsilon}(v)=\varepsilon(u)$, on a donc
% block display | line 11
\[
g(\gamma)=\gamma(v^{-1}u)-v^{-1}u\qquad\text{d'où }g\sim0\text{ dans }H^1(\Lambda,A)
\]
% block text | lines 12-17
(additif !). Comme $\delta=[\Gamma:\Lambda]$ est tel que $x\mapsto\delta x$ dans $A$ est \uline{bijectif} (car il l'est dans les $m^iA/m^{i+1}A$, qui sont des espaces vectoriels sur $k$) \deleted{\unsure{on a}} $H^1(\Gamma,A)\mathbin{-}H^1(\Lambda,A)$ est \uline{injectif}, donc il existe $w\in A$ tel que
% block display | line 18
\[
g(\gamma)=\gamma(w)-w\qquad\text{pour }\underline{\forall}\gamma\in\Gamma.
\]
% block text | lines 18-22
Donc \deleted{$\unsure{1+g(\gamma)=wg|1+g(\gamma)=w\gamma}$} On peut (ajoutant au besoin \deleted{\unsure{un}} un élément de $A^\Gamma$) supposer $w\in A^*$, \deleted{\unsure{Soit donc}} \inserted{et} $=\varepsilon(w)=1$. Alors on a \deleted{\unsure{donc}} \unsure{on a}
% block display | lines 23-25
\[
\begin{aligned}
g(\gamma)&=\varepsilon(w)g(\gamma)=wg(\gamma),\qquad\text{d'où}\\
g(\gamma)&=w^{-1}\gamma(w)-1\qquad\text{d'où}\\
1+g(\gamma)&=w^{-1}\gamma(w)
\end{aligned}
\]
% block text | line 26
donc $1+g(\gamma)\sim0$ en $H^1(\Gamma,A^*)$.
